\subsection{\q{Block out}ゲームのハイスコアファイルと原始的なシリアライゼーション}

多くのビデオゲームにはハイスコアファイルがあり、時々\q{Hall of fame}と呼ばれます。
古い\q{Block out}\footnote{\url{http://www.bestoldgames.net/eng/old-games/blockout.php}}ゲーム
（1989年の3Dテトリス）も例外ではありません。最後に見ることができるのは次のとおりです：

\begin{figure}[H]
\centering
\includegraphics[width=0.7\textwidth]{advanced/550_more_structs/blockout/hs.png}
\caption{ハイスコアテーブル}
\end{figure}

名前を追加した後にファイルが変更されたことがわかります。そのファイルは\IT{BLSCORE.DAT}です。
\begin{lstlisting}
% xxd -g 1 BLSCORE.DAT

00000000: 0a 00 58 65 6e 69 61 2e 2e 2e 2e 2e 00 df 01 00  ..Xenia.........
00000010: 00 30 33 2d 32 37 2d 32 30 31 38 00 50 61 75 6c  .03-27-2018.Paul
00000020: 2e 2e 2e 2e 2e 2e 00 61 01 00 00 30 33 2d 32 37  .......a...03-27
00000030: 2d 32 30 31 38 00 4a 6f 68 6e 2e 2e 2e 2e 2e 2e  -2018.John......
00000040: 00 46 01 00 00 30 33 2d 32 37 2d 32 30 31 38 00  .F...03-27-2018.
00000050: 4a 61 6d 65 73 2e 2e 2e 2e 2e 00 44 01 00 00 30  James......D...0
00000060: 33 2d 32 37 2d 32 30 31 38 00 43 68 61 72 6c 69  3-27-2018.Charli
00000070: 65 2e 2e 2e 00 ea 00 00 00 30 33 2d 32 37 2d 32  e........03-27-2
00000080: 30 31 38 00 4d 69 6b 65 2e 2e 2e 2e 2e 2e 00 b5  018.Mike........
00000090: 00 00 00 30 33 2d 32 37 2d 32 30 31 38 00 50 68  ...03-27-2018.Ph
000000a0: 69 6c 2e 2e 2e 2e 2e 2e 00 ac 00 00 00 30 33 2d  il...........03-
000000b0: 32 37 2d 32 30 31 38 00 4d 61 72 79 2e 2e 2e 2e  27-2018.Mary....
000000c0: 2e 2e 00 7b 00 00 00 30 33 2d 32 37 2d 32 30 31  ...{...03-27-201
000000d0: 38 00 54 6f 6d 2e 2e 2e 2e 2e 2e 2e 00 77 00 00  8.Tom........w..
000000e0: 00 30 33 2d 32 37 2d 32 30 31 38 00 42 6f 62 2e  .03-27-2018.Bob.
000000f0: 2e 2e 2e 2e 2e 2e 00 77 00 00 00 30 33 2d 32 37  .......w...03-27
00000100: 2d 32 30 31 38 00                                -2018.
\end{lstlisting}

すべてのエントリが明確に見えます。
最初のバイトはおそらくエントリの数です。
2番目はゼロで、実際、エントリの数は最初の2バイトにまたがる16ビット値である可能性があります。

次に、\q{Xenia}の名前の後に0xDFと0x01バイトが見えます。
Xeniaのスコアは479で、これは16進数で0x1DFです。
そのため、ハイスコア値は16ビット整数、またはおそらく32ビット整数でしょう：その後に2つのゼロバイトがあります。

配列要素と構造体要素の両方が常にメモリ内で隣接して配置されるという事実について考えてみましょう。
\myindex{\CStandardLibrary!write()}
\myindex{\CStandardLibrary!fwrite()}
\myindex{\CStandardLibrary!read()}
\myindex{\CStandardLibrary!fread()}
これにより、単純な\IT{write()}または\IT{fwrite()}関数を使用して配列/構造体全体をファイルに書き込み、
その後\IT{read()}または\IT{fread()}を使用してそれを復元することができます。これは非常に簡単です。
これが現在\IT{シリアライゼーション}と呼ばれているものです。

\subsubsection{読み込み}

ハイスコアファイルを読み込むCプログラムを書いてみましょう：

\begin{lstlisting}[style=customc]
#include <assert.h>
#include <stdio.h>
#include <stdint.h>
#include <string.h>

struct entry
{
	char name[11]; // incl. terminating zero
	uint32_t score;
	char date[11]; // incl. terminating zero
} __attribute__ ((aligned (1),packed));

struct highscore_file
{
	uint8_t count;
	uint8_t unknown;
	struct entry entries[10];
} __attribute__ ((aligned (1), packed));

struct highscore_file file;

int main(int argc, char* argv[])
{
	FILE* f=fopen(argv[1], "rb");
	assert (f!=NULL);
	size_t got=fread(&file, 1, sizeof(struct highscore_file), f);
	assert (got==sizeof(struct highscore_file));
	fclose(f);
	for (int i=0; i<file.count; i++)
	{
		printf ("name=%s score=%d date=%s\n",
				file.entries[i].name,
				file.entries[i].score,
				file.entries[i].date);
	};
};
\end{lstlisting}

すべての構造体フィールドが1バイト境界でパックされるように、GCC \IT{((aligned (1),packed))}属性が必要です。

もちろん、動作します：

\begin{lstlisting}
name=Xenia..... score=479 date=03-27-2018
name=Paul...... score=353 date=03-27-2018
name=John...... score=326 date=03-27-2018
name=James..... score=324 date=03-27-2018
name=Charlie... score=234 date=03-27-2018
name=Mike...... score=181 date=03-27-2018
name=Phil...... score=172 date=03-27-2018
name=Mary...... score=123 date=03-27-2018
name=Tom....... score=119 date=03-27-2018
name=Bob....... score=119 date=03-27-2018
\end{lstlisting}

（言うまでもなく、各名前は画面上とファイル内の両方でドットでパディングされています。おそらく美的な理由でしょう。）

\subsubsection{書き込み}

スコア値の幅について正しいかどうか確認しましょう。本当に32ビットでしょうか？

\begin{lstlisting}[style=customc]
int main(int argc, char* argv[])
{
	FILE* f=fopen(argv[1], "rb");
	assert (f!=NULL);
	size_t got=fread(&file, 1, sizeof(struct highscore_file), f);
	assert (got==sizeof(struct highscore_file));
	fclose(f);

	strcpy (file.entries[1].name, "Mallory...");
	file.entries[1].score=12345678;
	strcpy (file.entries[1].date, "08-12-2016");

	f=fopen(argv[1], "wb");
	assert (f!=NULL);
	got=fwrite(&file, 1, sizeof(struct highscore_file), f);
	assert (got==sizeof(struct highscore_file));
	fclose(f);
};
\end{lstlisting}

Blockoutを実行してみましょう：

\begin{figure}[H]
\centering
\includegraphics[width=0.7\textwidth]{advanced/550_more_structs/blockout/hs345678.png}
\caption{ハイスコアテーブル}
\end{figure}

最初の2桁（1と2）は切り捨てられています：12345678は345678になります。おそらく、これは書式の問題でしょう...しかし数値はほぼ正しいです。
今度は999999に変更して再度実行します：

\begin{figure}[H]
\centering
\includegraphics[width=0.7\textwidth]{advanced/550_more_structs/blockout/hs999999.png}
\caption{ハイスコアテーブル}
\end{figure}

今度は正しいです。はい、ハイスコア値は32ビット整数です。

\subsubsection{これはシリアライゼーションでしょうか？}

\dots{}ほとんどそうです。
このようなシリアライゼーションは、効率と速度が\ac{XML}や\ac{JSON}への変換や逆変換よりもはるかに重要な科学・エンジニアリングソフトウェアで非常に人気があります。

重要なことの一つは、ファイルをメモリにロードするたびにすべての構造体が異なる場所に割り当てられる可能性があるため、明らかにポインタをシリアライズできないということです。

しかし：ある種の低コストの\ac{MCU}上で、単純な\ac{OS}上で作業し、構造体が常にメモリ内の同じ場所に割り当てられている場合、おそらくポインタも保存・復元できるでしょう。

\subsubsection{ランダムノイズ}

この例を準備する際、\q{Block out}を何度も実行し、少し遊んでハイスコアテーブルをランダムな名前で埋める必要がありました。

ファイルにエントリが3つだけあった時、次のようなものを見ました：

\begin{lstlisting}
00000000: 03 00 54 6f 6d 61 73 2e 2e 2e 2e 2e 00 da 2a 00  ..Tomas.......*.
00000010: 00 30 38 2d 31 32 2d 32 30 31 36 00 43 68 61 72  .08-12-2016.Char
00000020: 6c 69 65 2e 2e 2e 00 8b 1e 00 00 30 38 2d 31 32  lie........08-12
00000030: 2d 32 30 31 36 00 4a 6f 68 6e 2e 2e 2e 2e 2e 2e  -2016.John......
00000040: 00 80 00 00 00 30 38 2d 31 32 2d 32 30 31 36 00  .....08-12-2016.
00000050: 00 00 57 c8 a2 01 06 01 ba f9 47 c7 05 00 f8 4f  ..W.......G....O
00000060: 06 01 06 01 a6 32 00 00 00 00 00 00 00 00 00 00  .....2..........
00000070: 00 00 00 00 00 00 00 00 00 00 00 00 00 00 00 00  ................
00000080: 00 00 00 00 00 00 00 00 00 00 00 00 00 00 00 00  ................
00000090: 00 00 00 00 00 00 00 00 00 00 00 00 00 00 00 00  ................
000000a0: 00 00 00 00 00 00 00 00 00 00 93 c6 a2 01 46 72  ..............Fr
000000b0: 8c f9 f6 c5 05 00 f8 4f 00 02 06 01 a6 32 06 01  .......O.....2..
000000c0: 00 00 98 f9 f2 c0 05 00 f8 4f 00 02 a6 32 a2 f9  .........O...2..
000000d0: 80 c1 a6 32 a6 32 f4 4f aa f9 39 c1 a6 32 06 01  ...2.2.O..9..2..
000000e0: b4 f9 2b c5 a6 32 e1 4f c7 c8 a2 01 82 72 c6 f9  ..+..2.O.....r..
000000f0: 30 c0 05 00 00 00 00 00 00 00 a6 32 d4 f9 76 2d  0..........2..v-
00000100: a6 32 00 00 00 00                                .2....
\end{lstlisting}

最初のバイトは3の値を持ち、3つのエントリがあることを意味します。
そして3つのエントリが存在します。
しかし、その後ファイルの後半にランダムノイズが見えます。

ノイズはおそらく初期化されていないデータに起源があります。
おそらく、\q{Block out}は\gls{heap}のどこかに10エントリ用のメモリを割り当て、そこには明らかに（他の何かからの）疑似ランダムノイズが存在していました。
その後、最初/2番目のバイトを設定し、3つのエントリを埋めましたが、残り7つのエントリには触れなかったため、それらはそのままファイルに書き込まれます。

\q{Block out}が次回の実行でハイスコアファイルをロードするとき、最初/2番目のバイトからエントリ数（3）を読み取り、その後は完全に無視します。

これは一般的な問題です。
厳密な意味での問題ではありません：バグではありませんが、情報が外部に漏洩する可能性があります。

\myindex{Microsoft Word}
1990年代のMicrosoft Wordバージョンは、以前に編集されたテキストの断片を*.doc*ファイルに残すことがよくありました。
当時、誰かから\IT{.doc}ファイルを取得し、16進エディタで開いて他の何かを読むのは、ある種の娯楽でした。
そのコンピュータで以前に編集されたものを読むことができました。

\myindex{Heartbleed}
\myindex{OpenSSL}
問題はさらに深刻になる可能性があります：OpenSSLのHeartbleedバグ\footnote{\url{https://en.wikipedia.org/wiki/Heartbleed}}。

\subsubsection{宿題}

\q{Block out}にはいくつかのポリキューブ（flat/basic/extended）があり、ピットのサイズを設定できるなどがあります。
そして、各設定について、\q{Block out}は独自のハイスコアテーブルを持っているようです。
一部の情報はおそらく\IT{BLSCORE.IDX}ファイルに格納されていることに気付きました。
これは、ハードコアな\q{Block out}ファンにとって宿題になる可能性があります---その構造も理解することです。

\q{Block out}ファイルはここにあります：\url{http://beginners.re/examples/blockout.zip}
（この例で使用したバイナリハイスコアファイルを含む）。
DosBoxを使用して実行できます。